% Generated draft: requires mathematical and source review.
\documentclass[11pt]{article}
\usepackage[margin=25mm]{geometry}
\usepackage{fontspec}
\setmainfont{Latin Modern Roman}
\usepackage{amsmath,amssymb,mathtools,unicode-math}
\setmathfont{Latin Modern Math}
\usepackage{enumitem,needspace,xurl,hyperref}
\setlength{\parindent}{0pt}
\setlength{\parskip}{5pt}
\setlength{\emergencystretch}{3em}
\newcommand{\N}{\mathbb{N}}
\newcommand{\Z}{\mathbb{Z}}
\newcommand{\Q}{\mathbb{Q}}
\newcommand{\R}{\mathbb{R}}
\newcommand{\C}{\mathbb{C}}
\newcommand{\F}{\mathbb{F}}
\newcommand{\A}{\mathbb{A}}
\newcommand{\PP}{\mathbb{P}}
\newcommand{\E}{\mathbb{E}}
\newcommand{\eps}{\varepsilon}
\newcommand{\dd}{\,\mathrm{d}}
\begin{document}
\section*{UnsolvedMath ID: 20001953}
\textbf{Source problem code:} catalogue entry
\subsection{Definitions and mathematical statement}
C. Boyer's Question 2.6 asks whether fixed sets of smaller subtori can yield cohomological information about abelian 3-Sasakian reductions. For a torus \(T\) preserving the 3-Sasakian structure on \(M\), write the contact moment map as
\[
 \mu:M\to\mathfrak t^*\otimes\mathbb R^3,\qquad
 \mu_i^\xi(p)=\eta_i(\xi_M(p)),\qquad Z=\mu^{-1}(0).
\]
The reduction is \(M\mathbin{/\!/\!/}T=Z/T\), with the usual admissibility assumptions. Cohomology in the source is rational. Its motivating sphere case has an \(S^1\)-action on \(S^{4n+3}\), locally free on the entire sphere and free on \(Z\), and reports surjectivity of
\[
 H^*_{S^1}(S^{4n+3};\mathbb Q)
 \longrightarrow H^*(S^{4n+3}\mathbin{/\!/\!/}S^1;\mathbb Q)
\]
``up to the middle dimension.'' Exact degree and reduction hypotheses are not supplied. The source notes that the fixed sets \(M^H\), \(H\subset T\), are contact, with spheres occurring among known Sasakian examples.

The aid gives an obstruction to a particular fixed-subtorus induction. For a positive-dimensional connected subtorus \(H\), the \(\mathfrak h\)-components of \(\mu\) vanish on \(M^H\). Thus its restriction is the residual \(T/H\) moment map, and
\[
 (\mu|_{M^H})^{-1}(0)=M^H\cap Z.
\]
If \(T\) is locally free on \(Z\), the right side is empty, because points there have finite stabilizers. Hence the residual zero reduction of every positive-dimensional fixed stratum is empty. This does not rule out every indirect use of fixed sets for cohomology; it rules out treating these as nonempty lower-rank zero reductions.

The aid separately reports that a linear torus action locally free on an entire sphere has rank at most one. A weight vector has stabilizer containing the kernel of its character, which has positive dimension if the torus rank exceeds one. This statement is explicitly about linear sphere actions, not arbitrary nonlinear torus actions. Finite isotropy subgroups are not positive-dimensional subtori and are not excluded by the fixed-stratum observation.

Under the stronger assumption that the action is locally free on all of \(M\), let \(X=M/T\) and \(Y=Z/T\) be the underlying coarse quotient pair. Rational equivariant cohomology identifies with ordinary rational cohomology of these quotients. The reported Kirwan cokernel description, in degree \(q\geq0\), is
\[
 \operatorname{coker}\bigl(\kappa_q:H_T^q(M;\mathbb Q)
                           \to H_T^q(Z;\mathbb Q)\bigr)
 \cong\ker\bigl(H^{q+1}(X,Y;\mathbb Q)
                         \to H^{q+1}(X;\mathbb Q)\bigr),
\]
using the long exact sequence of the pair. It is a reduction to a relative-cohomology calculation, not an assertion that the kernel vanishes. Local freeness only on \(Z\) does not justify replacing \(H_T^*(M)\) by \(H^*(M/T)\).

The concatenated Comment 2.7, also by C. Boyer, names [BG99, BGM94, BGMR98] as the main references for toral reductions. The trailing heading ``3 Kirwan surjectivity for contact quotients'' contains no further question in this record.
\subsection{Short English statement}
The supplied aid shows that positive-dimensional fixed subtori miss a locally free full zero level, so their residual zero reductions cannot directly support the proposed induction. Under global local freeness, a quotient-pair exact sequence describes the rational Kirwan cokernel. The general request for cohomological information and other possible uses of fixed sets remain distinct.
\subsection{Sources}
\begin{itemize}
\item Local source record: catalogue entry, catalogue entry; canonical \texttt{aim-geometry-notes.json}, \texttt{notes[290]}, source item 2.6. Title: ``Why fixed subtori miss the regular 3-Sasakian zero set.''
\item AIM August 2004 workshop list \emph{Moment maps and surjectivity in various geometries}, Question 2.6 and Comment 2.7 attributed to C. Boyer; supplied reference \url{https://aimath.org/WWN/momentmaps/momentmaps.pdf}. This draft uses the supplied text without accessing a PDF.
\item Original Comment 2.7 supplies [BG99, BGM94, BGMR98] for 3-Sasakian toral reductions. Full bibliographic entries are absent. The fixed-stratum obstruction, linear-sphere rank observation and cokernel formula are claims of research attempt 1.
\item Registry dated 14 August 2026: ``partially\_solved,'' ``PARTIAL-PROGRESS,'' statement status ``exact''; attempt partial, research-run status open, rigor ``proved,'' novelty confidence low. The machine-generated, non-peer-reviewed aid requires independent verification before citation.
\end{itemize}
\end{document}
